\section{The actual nonmonomial residual along general NC separation}
\label{sec:hironaka-normal-crossings-actual-residual}

Throughout this section the ambient complex dimension is positive.

Let $H$ be a genuine normal-crossings carrier, allowing local branches and
self-intersections, and let $F$ be an actual finite indexed weighted SNC
boundary adapted to $H$. Put $P=\prod_i\mathcal I_{F_i}^{a_i}$ and
$I=P\mathcal I_H$. The entire boundary and every zero or empty historical
entry are retained. The residual below is the actual intrinsic colon
residual, rather than a supplied factorization datum. This is the first
entry of the general marked invariant; equality of its lower history and
colexicographic entries is a separate question.

\begin{lemma}[The whole intrinsic monomial factor of the actual joint ideal]
\label{hir:general-nc-actual-joint-monomial-factor}
\uses{hir:general-nc-original-joint-pullback-exponents,hir:nc-ordered-tail-gamma-singleton-formula,hir:intrinsic-supported-ideal-monomial}
\lean{Hironaka.HironakaNormalCrossingsActualResidual.monomialPart_joint_eq_componentPowerProduct}
\leanok
If the ambient dimension is positive, then the whole intrinsic monomial
ideal of $I$ on the literal boundary $F$ equals $P$.
\end{lemma}
\begin{proof}
The actual joint coordinate-product theorem computes the intrinsic
exponent of $I$ at every boundary member through every point as its
literal weight $a_i$. The actual ordinary SNC product theorem computes
the exponent of $P$ on that same member as the same weight. Thus the
definitions of the two intrinsic monomial ideals give equality at every
stalk. Each factor ideal of $P$ is invertible, hence so is their actual
finite product. The actual marking-one cosupport formula for that
product says that every support point lies on a positive-weight member,
so its support lies in the union of the literal boundary. The proved
SNC-supported divisor theorem identifies the intrinsic monomial ideal
of this same whole product with $P$. Combining the two equalities proves
the assertion. Empty members and zero weights contribute no support.
\end{proof}

\begin{theorem}[The whole residual is the actual strict original ideal]
\label{hir:general-nc-actual-joint-residual}
\uses{hir:general-nc-actual-joint-monomial-factor,hir:monomial-part-ambient,hir:general-nc-whole-original-ideal-invertible,hir:monomial-part-along}
\lean{Hironaka.HironakaNormalCrossingsActualResidual.nonmonomialPart_joint_eq_original,Hironaka.HironakaNormalCrossingsActualResidual.nonmonomialPartAlong_joint_eq_original}
\leanok
In positive ambient dimension the entire ambient nonmonomial part of
$I$ on $F$ is $\mathcal I_H$. The same equality holds for the actual
nonmonomial part along the whole ambient manifold.
\end{theorem}
\begin{proof}
The whole NC ideal and the finite boundary product are invertible, so
their product $I$ is invertible. The genuine intrinsic factorization is
$I=\mathcal N(I)\mathcal M(I)$. Its residual is invertible, by the actual
colon theorem for nested invertible ideals. Substituting the preceding
whole monomial equality and commuting factors gives
$P\mathcal N(I)=P\mathcal I_H$. Cancellation of the actual invertible
factor $P$ gives the ambient equality. Along the whole manifold, a
restricted span is exactly its ambient span: apply the defining
neighbourhood equality at every point of the whole space. Thus the
ambient extension of a coherent ideal along the whole manifold is the
same ideal. The boundary exponents along the whole manifold are the
ambient exponents, and the power-ideal spans give the same membership.
Consequently both actual monomial ideals and both actual colon
residuals agree, proving the second assertion.
\end{proof}

\begin{theorem}[The genuine first invariant entry at every actual prefix]
\label{hir:general-nc-actual-prefix-first-entry}
\uses{hir:general-nc-actual-joint-residual,hir:general-nc-generated-exact-original-weights,hir:general-nc-original-joint-pullback-exponents,bcalc:nonmonomial-order-local}
\lean{Hironaka.HironakaNormalCrossingsActualResidual.higherSeparationCentres_nonmonomialPartAlong,Hironaka.HironakaNormalCrossingsActualResidual.higherSeparationCentres_nmOrd}
\leanok
At every actual recorded higher-stratum separation prefix, the entire
nonmonomial ideal of the same original marking-one controlled transform
is the vanishing ideal of its genuine differential strict original.
Its actual nonmonomial order at every point is the intrinsic number of
branches of that strict original.
\end{theorem}
\begin{proof}
The actual higher-centre induction produces the literal current
boundary-adapted strict original and the whole controlled-ideal
dictionary $I_T=P_T\mathcal I_{H_T}$, where $P_T$ uses precisely the
recorded historical weights $r-1$. The recorded carrier list is exactly
the current marked boundary, and every boundary factor is genuinely
invertible. Apply the preceding whole residual theorem to this actual
pair and product, and transport across those literal equalities. For
each natural $q$, membership in its actual residual order-$q$ set is
therefore equivalent to $q$ being at most the intrinsic branch count,
by the exact whole original-order theorem. The count plus one does not
belong to the next order set, whereas each positive integer up to the
count does. The definition as the least failed successor order thus
gives equality of the natural nonmonomial order with that count.
\end{proof}

\begin{corollary}[Every genuine separation centre is the whole first-entry maximum]
\label{hir:general-nc-actual-first-entry-maximum-centre}
\uses{hir:general-nc-actual-prefix-first-entry,hir:general-nc-internal-original-test-tower,hir:general-nc-bounded-actual-separation-tower}
\lean{Hironaka.HironakaNormalCrossingsActualResidual.higherSeparationCentre_eq_nmOrd_maxLocus}
\leanok
At every actual higher-stratum separation prefix, its next genuine
highest-branch centre is exactly the whole maximum locus of the actual
nonmonomial-order function on the same current marked cosupport.
\end{corollary}
\begin{proof}
Let $r\ge2$ be the attained actual global branch maximum defining that
centre. The preceding equality bounds the actual nonmonomial order of
every point by $r$. Every centre point belongs to the current strict
original, hence to the same current controlled cosupport, and has
order $r$; it therefore belongs to the asserted maximum locus.
Conversely take a point of that maximum locus. Choose an actual point
of the nonempty geometric centre. This chosen point lies in the same
cosupport and has order $r$, so maximality forces the first point to
have order at least $r$. The first-entry identity puts it in the
geometric highest-branch centre. This proves both whole-set inclusions.
\end{proof}

\begin{remark}
These are whole residual and first-entry identities for the genuine
original NC ideal, at the actual generated prefixes. They do not
identify all lower history-sensitive entries of the general canonical
invariant, or assert that its full maximum locus equals every geometric
highest-branch locus. An unrelated additional marked ideal can have
cosupport disjoint from the NC singular locus, so admissibility of the
geometric separation for that ideal requires a different construction.
The full canonical marked-ideal G9 endpoint remains open.
\end{remark}
